% --- Chapter: Functions and Closures ---
\section{Functions and Closures}

\subsection{Named Functions}

CRISP supports named function definition with the \texttt{fn} keyword. Functions are first-class values stored in the environment, enabling recursion and self-reference:

\begin{tcolorbox}[colback=codebg, colframe=darkpurple, colbacktitle=darkpurple, title=\textbf{\textcolor{codegold}{Named Functions}}, fonttitle=\bfseries, sharp corners]
\begin{lstlisting}
# Basic function definition
fn greet(name) {
    say "Hello, ", name, "!";
}

greet("CRISP");  # Hello, CRISP!

# Function with return value
fn add(a, b) {
    return a + b;
}

let result = add(3, 7);
say result;  # 10

# Early return
fn abs(x) {
    if x >= 0 {
        return x;
    }
    return -x;
}

say abs(-42);  # 42
say abs(42);   # 42
\end{lstlisting}
\end{tcolorbox}

\subsection{Function Parameters}

Functions accept positional parameters. Each parameter is bound to the corresponding argument at call time:

\begin{tcolorbox}[colback=codebg, colframe=darkpurple, colbacktitle=darkpurple, title=\textbf{\textcolor{codegold}{Function Parameters}}, fonttitle=\bfseries, sharp corners]
\begin{lstlisting}
# Multiple parameters
fn distance(x1, y1, x2, y2) {
    let dx = x2 - x1;
    let dy = y2 - y1;
    return sqrt(dx * dx + dy * dy);
}

# Parameters with sigils (type annotation)
fn process(@arr, %opts) {
    say "Array: ", @arr;
    say "Options: ", %opts;
}

process([1, 2, 3], { debug => true });
\end{lstlisting}
\end{tcolorbox}

\subsection{Recursion}

Functions are recursive with a maximum depth of 1000 calls. The function name is available inside its own body via self-reference:

\begin{tcolorbox}[colback=codebg, colframe=darkpurple, colbacktitle=darkpurple, title=\textbf{\textcolor{codegold}{Recursion}}, fonttitle=\bfseries, sharp corners]
\begin{lstlisting}
# Factorial — classic recursion
fn factorial(n) {
    if n <= 1 {
        return 1;
    }
    return n * factorial(n - 1);
}

say factorial(5);   # 120
say factorial(10);  # 3628800

# Fibonacci
fn fib(n) {
    if n <= 1 { return n; }
    return fib(n - 1) + fib(n - 2);
}

say fib(10);  # 55

# Mutual recursion
fn is_even(n) {
    if n == 0 { return true; }
    return is_odd(n - 1);
}

fn is_odd(n) {
    if n == 0 { return false; }
    return is_even(n - 1);
}
\end{lstlisting}
\end{tcolorbox}

\subsection{Lambda Functions}

Lambdas are anonymous functions defined with pipe syntax. They come in two forms: expression-body (implicit return) and block-body (explicit return):

\begin{tcolorbox}[colback=codebg, colframe=darkpurple, colbacktitle=darkpurple, title=\textbf{\textcolor{codegold}{Lambda Functions}}, fonttitle=\bfseries, sharp corners]
\begin{lstlisting}
# Expression-body lambda (implicit return)
let double = |x| => x * 2;
say double(10);  # 20

# Block-body lambda (explicit return)
let square = |x| => {
    let result = x * x;
    return result;
};
say square(5);  # 25

# Multiple parameters
let add = |a, b| => a + b;
say add(3, 7);  # 10

# Lambda as argument
let numbers = [1, 2, 3, 4, 5];
let squared = map(numbers, |n| => n * n);
say squared;  # [1, 4, 9, 16, 25]

# Lambda capturing outer variables
let factor = 10;
let scale = |x| => x * factor;
say scale(5);  # 50
\end{lstlisting}
\end{tcolorbox}

\subsection{Closures}

Functions capture their lexical environment at definition time. Each function call creates a fresh environment that chains to the captured parent:

\begin{tcolorbox}[colback=codebg, colframe=darkpurple, colbacktitle=darkpurple, title=\textbf{\textcolor{codegold}{Closures}}, fonttitle=\bfseries, sharp corners]
\begin{lstlisting}
# Counter factory — returns a closure
fn make_counter(start) {
    return || => {
        let current = start;
        start = start + 1;
        return current;
    };
}

let counter1 = make_counter(0);
let counter2 = make_counter(100);

say counter1();  # 0
say counter1();  # 1
say counter1();  # 2
say counter2();  # 100
say counter2();  # 101

# Each closure has its own captured environment
\end{lstlisting}
\end{tcolorbox}

\subsection{Functions as Values}

Functions are first-class values: they can be stored in variables, passed as arguments, and returned from other functions:

\begin{tcolorbox}[colback=codebg, colframe=darkpurple, colbacktitle=darkpurple, title=\textbf{\textcolor{codegold}{First-Class Functions}}, fonttitle=\bfseries, sharp corners]
\begin{lstlisting}
# Assign function to variable
let greet = fn(name) {
    return "Hello, " + name + "!";
};
say greet("CRISP");  # Hello, CRISP!

# Pass function as argument
fn apply_twice(f, x) {
    return f(f(x));
}

let inc = |x| => x + 1;
say apply_twice(inc, 5);  # 7 (inc(inc(5)) = inc(6) = 7)

# Return function from function
fn multiplier(factor) {
    return |x| => x * factor;
}

let triple = multiplier(3);
say triple(10);  # 30
\end{lstlisting}
\end{tcolorbox}
